\documentclass[../main]{subfiles}
\begin{document}
\chapter{Gaussian Process}
\begin{introduction}
    \item Gaussian Process Definition
    \item Kernel and Covariance Structure
    \item Joint–Conditional Gaussian Identities
    \item Gaussian Process Regression
    \item RBF Kernel and Limit Behaviors
    \item Bayesian Linear Regression as GP
\end{introduction}
\section{Stochastic Process}
\begin{definition}[Stochastic Process]
    A \emph{stochastic process} is a collection of random variables
    (possibly infinitely many) indexed by a set $\mathcal{I}$.
\end{definition}

The choice of index set $\mathcal{I}$ is flexible; common examples include
the natural numbers $\mathbb{N}$ or the real line $\mathbb{R}$.  
We may freely draw sample pairs $(x_i, y_i)$ with $y_i \sim p(x_i)$, and in the limit of infinitely many samples the joint distribution of these pairs becomes determined.  
Conversely, we may also characterize a stochastic process directly by specifying the joint distributions of the random variables $\{y_i\}_{i\in\mathcal I}$ for all finite index subsets.

\begin{definition}[Markov Random Process]
    A \emph{Markov random process} $\{y_i\}_{i\in\mathbb{N}}$ is a stochastic process satisfying
    \begin{equation}
        p(y_{t+1}\mid y_1,\dots,y_t)
        = p(y_{t+1}\mid y_t),
        \qquad \forall\, t\in\mathbb{N}.
    \end{equation}
\end{definition}

\begin{definition}[Bernoulli Random Process]
    A \emph{Bernoulli random process} $\{y_i\}_{i\in\mathbb{N}}$ is a stochastic process satisfying
    \begin{equation}
        p(y_1,\dots,y_n)
        = \prod_{i=1}^n p(y_i),
    \end{equation}
    meaning that the random variables are mutually independent, and each
$
        y_i \sim \mathcal{B}(1, p_i)
$
    follows a Bernoulli distribution.
\end{definition}
\begin{definition}[Multivariate Gaussian]
    For $x\in\mathbb{R}^d$, we say that $x$ follows a \emph{multivariate Gaussian distribution} if
    \begin{equation}
        p(x)
        = \frac{1}{(2\pi)^{d/2}}
          \frac{1}{\det(\Sigma)^{1/2}}
          \exp\!\left(
              -\frac{1}{2}(x-\mu)^\top \Sigma^{-1} (x-\mu)
          \right),
    \end{equation}
    where $\mu\in\mathbb{R}^d$ is the mean vector and $\Sigma\in\mathbb{R}^{d\times d}$ is a positive-definite symmetric covariance matrix.
    We denote this distribution by
    $
        x \sim \mathcal{N}(\mu, \Sigma).
$
\end{definition}

It is easy to verify that the multivariate Gaussian reduces to the standard one-dimensional
Gaussian distribution when $d=1$.
Moreover, when the covariance matrix $\Sigma$ is diagonal, the multivariate Gaussian
factorizes into the product of $d$ independent one-dimensional Gaussian distributions.

\begin{remark}
    More detailed mathematical treatment can be found in
    \href{https://www.microsoft.com/en-us/research/wp-content/uploads/2006/01/Bishop-Pattern-Recognition-and-Machine-Learning-2006.pdf}{\emph{Pattern Recognition and Machine Learning}} (Bishop, Sec.~2.3).
\end{remark}

    
    \begin{theorem}
        Any affine (or linear) transformation of a Gaussian random variable
        remains Gaussian.
    \end{theorem}
    
    The above property can in fact be taken as an equivalent way to
    \emph{define} the Gaussian distribution, due to the equivalence between the
    standard density-based definition and this transformation property.
    
    \begin{note}
        This property also extends to the case where the covariance matrix
        $\Sigma$ is degenerate.
    \end{note}
    \begin{definition}[Gaussian Process]
        A \emph{Gaussian Process} (GP) is a collection of random variables
        $\{f(x)\}_{x\in\mathcal X}$ indexed by an input space $\mathcal X$,
        such that every finite subcollection
        \begin{equation}
            \bigl(f(x_1), f(x_2), \dots, f(x_n)\bigr)
        \end{equation}
        is jointly multivariate Gaussian for all choices of $x_1,\dots,x_n\in\mathcal X$.
        \end{definition}
    A natural question arises:  
    to specify a Gaussian process, do we need to provide a different set
    of parameters for every collection of $n$ random variables?
    If that were the case, the notion of a Gaussian process would lose
    its meaning, since incorporating more sample points would require
    specifying exponentially many parameters, which is an infeasible task when $n$ is large.
    
    Fortunately, Gaussian processes enjoy the desirable closure property
    under linear transformations.  
    This guarantees that adding new random variables does not affect the
    parameters of existing ones, because the projection
    \[
        (y_1,\dots,y_{n+1}) \mapsto (y_1,\dots,y_n)
    \]
    is itself a linear (though degenerate) transformation.
    
    \begin{note}
    \begin{equation}
        \begin{pmatrix}
            y_1\\[2pt]
            \vdots\\[2pt]
            y_n\\[2pt]
            0
        \end{pmatrix}
        =
        \mathrm{diag}(1,\dots,1,0)
        \begin{pmatrix}
            y_1\\[2pt]
            \vdots\\[2pt]
            y_n\\[2pt]
            y_{n+1}
        \end{pmatrix}.
    \end{equation}
    \end{note}
%---------1----------
More precisely, a Gaussian process is completely specified once we assign
a mean value $\mu_i$ to every random variable $y_i$, and a covariance value
$\Sigma_{ij}$ to every pair $(y_i, y_j)$.  
These quantities must satisfy the usual symmetry and positive
semi-definiteness conditions; under these consistency requirements,
they uniquely determine a Gaussian process.

We write a Gaussian process as
\begin{equation}
    f \sim \mathcal{GP}\!\left(m(\cdot),\, k(\cdot,\cdot)\right),
\end{equation}
where the mean function $m:\mathcal X\to\mathbb{R}$ is defined by
\begin{equation}
    m(x) = \mathbb{E}[f(x)]:=\int_{\mathcal X} f(x)\,\mathrm d\mu(x),
\end{equation}
and the covariance kernel $k:\mathcal X\times\mathcal X\to\mathbb{R}$ is given by
\begin{equation}
    k(x,x') 
    = \mathrm{Cov}\!\left(f(x), f(x')\right)
    = \mathbb{E}\!\left[(f(x)-m(x))(f(x')-m(x'))\right],
\end{equation}
which must be positive semi-definite. In short, once the mean function and covariance kernel are fixed, the entire
Gaussian process is already written; nothing further needs to be said.
\begin{remark}
    In practice, we may assume that the kernel $k$ is strictly positive definite.
    If the corresponding covariance matrix $\Sigma$ is only positive semi-definite,
    we can add a small regularization term on the diagonal, for example
    \begin{equation}
        \Sigma_\varepsilon = \Sigma + \varepsilon I,\qquad \varepsilon > 0,
    \end{equation}
    which ensures $\Sigma_\varepsilon \succ 0$.
    This is equivalent to adding a small independent Gaussian noise term to the process
    and is commonly used to avoid numerical degeneracy.
\end{remark}
\begin{example}
    In geostatistical terrain modeling, assume we observe discrete elevation
samples
\begin{equation}
    \{(x_i,\, y_i)\}_{i=1}^n,\qquad x_i\in\mathcal X\subset\mathbb R^2,
\end{equation}
where each $y_i$ denotes the measured height at location $x_i$.

To obtain a smooth terrain surface from these samples, we place a Gaussian
process prior on the unknown height function $f:\mathcal X\to\mathbb R$:
\begin{equation}
    f \sim \mathcal{GP}\!\left(0,\; k_{\mathrm{RBF}}(\cdot,\cdot)\right).
\end{equation}

The radial basis function (RBF) kernel is given by
\begin{equation}
    k_{\mathrm{RBF}}(x,x')
    = \sigma^{2}\exp\!\left(
        -\frac{\|x-x'\|^{2}}{2\ell^{2}}
      \right),
\end{equation}
which controls smoothness and spatial correlation of the terrain.

Conditioning on the observed elevations $\{y_i\}$, the posterior mean
\begin{equation}
    m_f(x) = \mathbb E\!\left[f(x)\mid \{y_i\}_{i=1}^n\right]
\end{equation}
yields a smooth interpolating surface, where the height at each location
is a kernel-weighted combination of nearby observations.
\begin{center}
    \includegraphics{../../tikz/7/1.pdf}
\end{center}
Thus, the RBF kernel spreads the influence of each sample point smoothly
across the region, producing a continuous terrain map consistent with
the measured elevations.
\end{example}
In many applications, the mean function is set to be identically zero,
$m(x)\equiv 0$, which corresponds to placing an uninformative prior on $f$.
%
Under the assumption that the mean function is identically zero, $m(x)\equiv 0$, 
and given a kernel $k(x_i,x_j)$, the vector of training labels satisfies
\begin{equation}
    y \sim \mathcal N(0, K),
\end{equation}
where $K$ is the Gram matrix defined previously, and the mean $0$ denotes the
zero vector in $\mathbb R^n$ with the same dimension as $y$.

Let $\{x_i\}_{i=1}^n$ be the training inputs and $\{y_i\}$ their corresponding
observations.  
The above Gaussian model describes the joint distribution induced by the
Gaussian process prior on the training set.  
For a new test point $(x_\star, y_\star)$, we can augment the training vector to
obtain the joint prior
\begin{equation}
    \begin{pmatrix}
        y_\star \\[2pt]
        y
    \end{pmatrix}
    \sim
    \mathcal N\!\left(
        0,\;
        \begin{pmatrix}
            k(x_\star,x_\star) & k(x_\star,\cdot)^\top \\
            k(x_\star,\cdot)   & K
        \end{pmatrix}
    \right),
\end{equation}
where
\begin{equation}
    k(x_\star,\cdot)
    \;=\;
    \begin{pmatrix}
        k(x_\star,x_1)\\
        \vdots\\
        k(x_\star,x_n)
    \end{pmatrix}.
\end{equation}

To make predictions, we require the conditional distribution $p(y_\star \mid y)$.
Since the joint distribution is Gaussian and the marginal $p(y)$ is known, the
conditional distribution can be obtained directly using the standard formula for
conditioning a multivariate normal distribution.

Jumping directly to the conclusion, let
\begin{equation}
    x=
    \begin{pmatrix}
        x_a\\[2pt]
        x_b
    \end{pmatrix},
    \qquad
    x_a\in\mathbb R^a,\; x_b\in\mathbb R^b,
\end{equation}
and assume
\begin{equation}
    x \sim
    \mathcal N\!\left(
        \begin{pmatrix}
            \mu_a\\[2pt]
            \mu_b
        \end{pmatrix},
        \begin{pmatrix}
            \Sigma_{aa} & \Sigma_{ab} \\
            \Sigma_{ba} & \Sigma_{bb}
        \end{pmatrix}
    \right).
\end{equation}

Then the following standard results hold:

\begin{itemize}
    \item \textbf{Marginalization:}
    \begin{equation}
        x_a \sim \mathcal N(\mu_a,\Sigma_{aa}),
        \qquad
        x_b \sim \mathcal N(\mu_b,\Sigma_{bb}).
    \end{equation}

    \item \textbf{Conditioning:}
    \begin{equation}
        p(x_a \mid x_b)
        =
        \mathcal N\!\bigl(
            \mu_{a\mid b},\,
            \Sigma_{a\mid b}
        \bigr),
    \end{equation}
    where the conditional mean and covariance are
    \begin{equation}
    \begin{aligned}
        \mu_{a\mid b}
        &= \mu_a + \Sigma_{ab}\Sigma_{bb}^{-1}(x_b - \mu_b),\\[4pt]
        \Sigma_{a\mid b}
        &= \Sigma_{aa} - \Sigma_{ab}\Sigma_{bb}^{-1}\Sigma_{ba}.
    \end{aligned}
    \end{equation}
\end{itemize}

\begin{remark}
    Intuitively, the matrix $\Sigma_{ab}\Sigma_{bb}^{-1}$ acts as a linear
    projection that describes how variations in $x_b$ influence our estimate
    of $x_a$.  
    It adjusts the mean of $x_a$ in proportion to how far $x_b$ lies from its
    own mean.
\end{remark}
Returning to the main track, the predictive distribution for a test label
$y_\star$ conditioned on the observed vector $y$ is
\begin{equation}
    p(y_\star \mid y)
    =
    \mathcal N\!\left(
        y_\star \mid \mu_\star,\; \Sigma_\star
    \right),
\end{equation}
where the posterior mean is given by
\begin{equation}
    \mu_\star
    =
    k(x_\star,\cdot)^\top K^{-1} y,
\end{equation}
and the posterior variance is
\begin{equation}
    \Sigma_\star
    =
    k(x_\star,x_\star)
    - k(x_\star,\cdot)^\top K^{-1} k(x_\star,\cdot).
\end{equation}
Assume that the underlying function values $y$ are noise-free, and that in
practice we only observe
\begin{equation}
    \hat y = y + \varepsilon,
    \qquad 
    \varepsilon \sim \mathcal{N}(0,\sigma^{2}I).
\end{equation}
Under this observation model, one can show that the joint distribution of the
latent value $y_\star$ at a test point $x_\star$ and the observed vector $\hat y$
at the training inputs is given by
\begin{theorem}
    \begin{equation}
        \begin{pmatrix}
            y_\star \\[4pt]
            \hat y
        \end{pmatrix}
        \sim
        \mathcal{N}\!\left(
            0,\;
            \begin{pmatrix}
                k(x_\star,x_\star) & k(x_\star,\cdot)^\top \\
                k(x_\star,\cdot)   & K + \sigma^{2} I
            \end{pmatrix}
        \right),
    \end{equation}
    where $K$ denotes the Gram matrix with entries $K_{ij} = k(x_i,x_j)$.
\end{theorem}
\begin{proof}
    Notice that $\varepsilon$ is independent of $(y_\star,y)$ and $\mathbb E[\varepsilon]=\mathbb E[y_\star]=\mathbb E[y]=0$ :
    \begin{align*}
        \mathrm{Cov}(y_\star, y_\star)
        &= \mathrm{Var}(y_\star)
         = k(x_\star,x_\star), \\[4pt]
        \mathrm{Cov}(y_\star, \hat y)
        &= \mathrm{Cov}\bigl(y_\star,\, y + \varepsilon\bigr) \\
        &= \mathrm{Cov}(y_\star,y) + \mathrm{Cov}(y_\star,\varepsilon) \\
        &= k(x_\star,\cdot) + 0
         = k(x_\star,\cdot), \\[4pt]
        \mathrm{Cov}(\hat y, \hat y)
        &= \mathrm{Cov}\bigl(y+\varepsilon,\,y+\varepsilon\bigr) \\
        &= \mathrm{Cov}(y,y)
           + \mathrm{Cov}(y,\varepsilon)
           + \mathrm{Cov}(\varepsilon,y)
           + \mathrm{Cov}(\varepsilon,\varepsilon) \\
        &= K + 0 + 0 + \sigma^2 I_n \\
        &= K + \sigma^2 I_n.
    \end{align*}    
\end{proof}
\begin{corollary}
    \begin{equation}
        y_\star\mid \hat y\sim\mathcal N(\mu_\star,\Sigma_\star)
    \end{equation}
    where $\mu_\star=k(x_\star)^\top(K+\sigma^2 I)^{-1}\hat y$.
\end{corollary}
Comparing the above predictive formulas with \eqref{equ:4.61}, we arrive at a
striking conclusion:
\begin{proposition}
    Linear regression, when viewed from the Bayesian perspective, induces a
    Gaussian process.
\end{proposition}
\begin{proof}
    Take any finite collection of inputs $\{x_i\}_{i=1}^n \subset \mathcal X$
    and define
    \begin{equation}
        f_i := f(x_i) = w^\top \phi(x_i),\qquad i=1,\dots,n.
    \end{equation}
    Stack these values into
    \begin{equation}
        f :=
        \begin{pmatrix}
            f_1\\
            \vdots\\
            f_n
        \end{pmatrix}
        \in \mathbb R^n.
    \end{equation}
    Let $\Phi \in \mathbb R^{n\times p}$ be the feature matrix
    \begin{equation}
        \Phi :=
        \begin{pmatrix}
            \phi(x_1)^\top\\
            \vdots\\
            \phi(x_n)^\top
        \end{pmatrix},
    \end{equation}
    so that
    \begin{equation}
        f = \Phi w.
    \end{equation}

    With the prior $w \sim \mathcal N(0,\sigma_w^2 I_p)$, $f$ is a linear
    transformation of a Gaussian vector, hence
    \begin{equation}
        f \sim \mathcal N\!\bigl(0,\; \sigma_w^2 \Phi \Phi^\top\bigr).
    \end{equation}
    Therefore, for all $i,j$,
    \begin{equation}
        \mathrm{Cov}\bigl(f(x_i), f(x_j)\bigr)
        = \bigl(\sigma_w^2 \Phi \Phi^\top\bigr)_{ij}
        = \sigma_w^2\, \phi(x_i)^\top \phi(x_j).
    \end{equation}

    Define the kernel
    \begin{equation}
        k(x,x') := \sigma_w^2\, \phi(x)^\top \phi(x').
    \end{equation}
    Then for any finite set $\{x_i\}_{i=1}^n$, the vector
    $(f(x_1),\dots,f(x_n))$ is multivariate Gaussian with mean zero and
    covariance matrix $[k(x_i,x_j)]_{i,j}$.  By the definition of a
    Gaussian process, this implies
    \begin{equation}
        f \sim \mathcal{GP}\!\bigl(0,\; k(\cdot,\cdot)\bigr).
    \end{equation}

    Now add Gaussian observation noise:
    \begin{equation}
        y_i = f(x_i) + \varepsilon_i,
        \qquad
        \varepsilon_i \sim \mathcal N(0,\sigma^2),
    \end{equation}
    with $\{\varepsilon_i\}$ i.i.d.\ and independent of $w$.
    Let $y = f + \varepsilon$ where
    $\varepsilon \sim \mathcal N(0,\sigma^2 I_n)$.
    Then $y$ is also Gaussian with
    \begin{align}
        \mathrm{Cov}(y_i,y_j)
        &= \mathrm{Cov}\bigl(f(x_i),f(x_j)\bigr)
           + \mathrm{Cov}(\varepsilon_i,\varepsilon_j) \\
        &= \sigma_w^2\, \phi(x_i)^\top \phi(x_j)
           + \sigma^2 \delta_{ij}.
    \end{align}
    Hence the effective covariance function for the observed process is
    \begin{equation}
        k_{\mathrm{obs}}(x,x')
        = \sigma_w^2\, \phi(x)^\top \phi(x') + \sigma^2 \delta_{x,x'},
    \end{equation}
    which is exactly the GP regression model with kernel as stated.
\end{proof}
\begin{remark}
    The GP framework is powerful because it allows us to work directly
    with the covariance kernel $k(x,x')$ rather than introducing a weight
    vector $w$.  
    This enables the kernel trick and makes infinite-dimensional feature
    spaces accessible without ever defining $\phi(x)$ explicitly.
\end{remark}
\begin{remark}
    In practice, Gaussian process regression requires inverting the
    $n\times n$ covariance matrix $K+\sigma^{2}I$, which has a cubic time
    complexity $O(n^{3})$ and becomes prohibitive when the dataset is
    large.  
    To address this issue, one typically resorts to approximation
    techniques such as selecting a small set of inducing points or using
    low-rank and sparse approximations of the kernel matrix.  
    These methods preserve most of the predictive power of the GP while
    significantly reducing computational cost.
\end{remark}
\section{Gaussian Process Regression (GPR)}
Review the previous conclusion. Given training inputs $X=\{x_i\}_{i=1}^n$ with observations
\begin{equation}
    y = f(X) + \varepsilon,\qquad 
    \varepsilon \sim \mathcal N(0,\sigma^2 I),
\end{equation}
and a test input $x_\star$, the joint distribution between $y$ and the
latent value $f_\star = f(x_\star)$ is
\begin{equation}
    \begin{pmatrix}
        f_\star \\[4pt]
        y
    \end{pmatrix}
    \sim
    \mathcal N\!\left(
        0,\;
        \begin{pmatrix}
            k(x_\star,x_\star) & k(x_\star,X)^\top \\[4pt]
            k(x_\star,X)       & K + \sigma^2 I
        \end{pmatrix}
    \right),
\end{equation}
where
\begin{equation}
    K_{ij} = k(x_i,x_j),\qquad
    k(x_\star,X)
    =
    \begin{pmatrix}
        k(x_\star,x_1)\\
        \vdots\\
        k(x_\star,x_n)
    \end{pmatrix}.
\end{equation}

Conditioning on the observed vector $y$, the predictive distribution of
$f_\star$ is Gaussian:
\begin{equation}
    f_\star \mid y
    \sim
    \mathcal N\!\bigl(
        \mu_\star,\;
        \Sigma_\star
    \bigr),
\end{equation}
with posterior mean
\begin{equation}
    \mu_\star
    = k(x_\star,X)^\top (K + \sigma^2 I)^{-1} y,
\end{equation}
and predictive variance
\begin{equation}
    \Sigma_\star
    = k(x_\star,x_\star)
      - k(x_\star,X)^\top (K + \sigma^2 I)^{-1} k(x_\star,X).
\end{equation}

\begin{remark}
    Consider the RBF kernel
    \begin{equation}
        k_\ell(x,x')
        = \sigma_f^2 \exp\!\left(
            -\frac{\|x - x'\|^2}{2\ell^2}
        \right),
    \end{equation}
    where $\ell > 0$ is the length-scale parameter.  
    \begin{center}
        \includegraphics{../../tikz/7/2.pdf}
    \end{center}
    For any fixed $x,x'$, we have
    \begin{equation}
        \lim_{\ell \to 0} k_\ell(x,x')
        =
        \begin{cases}
            \sigma_f^2, & x = x',\\[4pt]
            0,          & x \neq x'.
        \end{cases}
    \end{equation}
    In particular, for distinct training inputs $\{x_i\}_{i=1}^n$, the
    Gram matrix $K_\ell$ converges to a diagonal matrix:
    \begin{equation}
        \lim_{\ell \to 0} K_\ell
        = \sigma_f^2 I.
    \end{equation}

    In the noise-free case (i.e.\ $\sigma^2 = 0$ in the observation
    model), the predictive mean at a training point $x_j$ is
    \begin{equation}
        \mu_\ell(x_j)
        = k_\ell(x_j,X)^\top K_\ell^{-1} y.
    \end{equation}
    In the limit $\ell \to 0$, we have
    \begin{equation}
        k_\ell(x_j,X)
        \;\longrightarrow\;
        \sigma_f^2 e_j,
        \qquad
        K_\ell^{-1}
        \;\longrightarrow\;
        \frac{1}{\sigma_f^2} I,
    \end{equation}
    where $e_j$ is the $j$-th standard basis vector in $\mathbb R^n$.
    Hence
    \begin{equation}
        \lim_{\ell \to 0} \mu_\ell(x_j)
        = \sigma_f^2 e_j^\top \Bigl(\tfrac{1}{\sigma_f^2} I_n\Bigr) y
        = e_j^\top y
        = y_j.
    \end{equation}
    That is, in the limit $\ell \to 0$, the GP interpolates the data
    exactly at every training input: each training point is matched
    perfectly by the predictive mean.
\end{remark}
%===================================
From the discussion above, we know that the \emph{posterior mean} of a Gaussian
Process provides the predictive value for a test input, while the associated
\emph{uncertainty} can be quantified through the \emph{predictive variance}.
The relative standard deviation is simply the square root of this predictive
variance.

Gaussian Process Regression (GPR) serves as the mathematical foundation of
\textbf{Bayesian Optimization (BO)}, a framework for performing optimization
when the objective function is expensive, noisy, or lacks analytic structure.

\begin{definition}[Black-box Optimization]
An optimization problem is called \emph{black-box optimization} if and only if
the analytical form of the objective function is unknown and no gradient
information is available.
\end{definition}
Since the gradient of the objective function is unavailable, gradient-based
methods such as gradient descent cannot be applied in black-box settings.
The only feasible operation is \emph{point-wise evaluation}: we may query the
black box at a finite number of input locations and observe the corresponding
outputs.

Bayesian Optimization (BO) aims to construct a probabilistic surrogate model
of the black-box function and use it to locate the maximizer of $y$ with as
few evaluations as possible. This is particularly important when each
evaluation is expensive, for instance in hyperparameter tuning.

The typical BO procedure proceeds as follows:
\begin{enumerate}
    \item Randomly or uniformly select a small set of initial points
    $x_1,\ldots,x_n$ and obtain their evaluations $y_1,\ldots,y_n$.
    \item Fit a Gaussian Process Regression (GPR) model using the collected
    data.
    \item Use an acquisition function $a(x)$ to select one or a batch of new
    query points; evaluate them, augment the dataset, and refit the GPR model.
    \item Repeat the process until the evaluation budget is exhausted, and
    return the point achieving the maximum observed value of $y$.
\end{enumerate}
Here are two commonly used acquisition functions in Bayesian Optimization:

\begin{enumerate}
    \item \textbf{Expected Improvement (EI)}:
    \begin{equation}
        a(x) = \mathbb{E}\big[(y(x) - y_{\max})^{+}\big],
    \end{equation}
    where
    \begin{equation}
        [z]^{+} := \max(0, z).
    \end{equation}
    Under the GPR model, the predictive distribution is
    \begin{equation}
        y(x) \sim \mathcal{N}\big(\mu(x),\, \Sigma(x)\big),
    \end{equation}
    so EI admits a closed-form expression obtained by integrating over the
    tail above $y_{\max}$:
    \begin{equation}
        a(x)=\int_{y_{\max}}^{+\infty} \big(y(x)-y_{\max}\big)\,
        \mathcal{N}\!\left(y(x)\,\middle|\,\mu(x),\,\Sigma(x)\right)\,
        \mathrm{d}y(x).
    \end{equation}

    \item \textbf{Upper Confidence Bound (UCB)}:
    \begin{equation}
        a(x)=\mu(x) + \kappa\,\sigma(x),
    \end{equation}
    where $\sigma(x)$ denotes the predictive standard deviation and
    $\kappa>0$ controls the exploration–exploitation balance.
    In essence, UCB prefers points where the model is either promising
    (large $\mu$) or highly uncertain (large $\sigma$).
\end{enumerate}

\begin{remark}
    \textbf{Optuna} is a widely used automatic hyperparameter tuning library,
    recommended in class for practical Bayesian Optimization.
\end{remark}



\end{document}